% Capitulo 3: los cuatro diagramas, cada uno con pie, etiqueta y referencia.
\section{Diagramas}\label{sec:diagramas}

texforge compone diagramas descritos como texto dentro del propio fuente,
sin Inkscape, sin Node y sin binarios externos: al compilar, cada entorno se
compone en Rust puro y el resultado entra al PDF como una figura más.
Los tres motores entienden las mismas opciones --- \texttt{width},
\texttt{pos}, \texttt{caption} y \texttt{label} --- y el estilo visual se
fija una vez en \texttt{project.toml} bajo \texttt{[diagrams]}, aquí el
preajuste \texttt{editorial}, de modo que los tres dibujos comparten tinta,
acento y aire aunque los dibujen motores distintos; la familia tipográfica
llega hasta donde cada motor la deja elegir.

\subsection{Mermaid}

La figura~\ref{fig:secuencia} muestra la secuencia de una compilación: el
usuario pide el PDF, texforge valida el fuente y, si hay diagramas, los
envía al motor de composición dentro de un bucle antes de producir el
resultado. El lenguaje es Mermaid~\cite{mermaid}, el de los diagramas de
secuencia y de flujo que ya se escriben en la documentación técnica: los
participantes se declaran con una palabra, los mensajes con flechas y el
texto de cada flecha viaja sobre ella sin sintaxis adicional.

\begin{mermaid}[width=0.9\linewidth, caption={Secuencia de compilación con texforge}, label={fig:secuencia}, pos=H]
  sequenceDiagram
  autonumber
  actor Usuario
  participant texforge
  participant Validador
  participant Motor
  Usuario->>texforge: texforge build
  activate texforge
  texforge->>Validador: analiza el fuente
  activate Validador
  Validador-->>texforge: fuente válido
  deactivate Validador
  alt hay diagramas
  texforge->>Motor: prepara el motor
  activate Motor
  loop cada diagrama
  texforge->>Motor: renderiza diagrama
  Motor-->>texforge: figura lista
  end
  Motor-->>texforge: motor libre
  deactivate Motor
  else sin diagramas
  Note over texforge,Motor: compila directo
  end
  texforge->>Motor: compone el PDF
  activate Motor
  Motor-->>texforge: PDF final
  deactivate Motor
  texforge-->>Usuario: build/ listo
  deactivate texforge
\end{mermaid}

\subsubsection{El mismo diagrama en los cuatro preajustes}

El atributo \texttt{style} aplica uno de los cuatro preajustes; la misma
secuencia se repite abajo, figuras \ref{fig:secuencia-editorial},
\ref{fig:secuencia-monochrome}, \ref{fig:secuencia-technical} y
\ref{fig:secuencia-default}.

\begin{mermaid}[style=editorial, width=0.62\linewidth, caption={\texttt{editorial}: paleta sobria con un solo acento}, label={fig:secuencia-editorial}, pos=t]
  sequenceDiagram
  autonumber
  actor Usuario
  participant texforge
  participant Validador
  participant Motor
  Usuario->>texforge: texforge build
  activate texforge
  texforge->>Validador: analiza el fuente
  activate Validador
  Validador-->>texforge: fuente válido
  deactivate Validador
  alt hay diagramas
  texforge->>Motor: prepara el motor
  activate Motor
  loop cada diagrama
  texforge->>Motor: renderiza diagrama
  Motor-->>texforge: figura lista
  end
  Motor-->>texforge: motor libre
  deactivate Motor
  else sin diagramas
  Note over texforge,Motor: compila directo
  end
  texforge->>Motor: compone el PDF
  activate Motor
  Motor-->>texforge: PDF final
  deactivate Motor
  texforge-->>Usuario: build/ listo
  deactivate texforge
\end{mermaid}

\begin{mermaid}[style=monochrome, width=0.62\linewidth, caption={\texttt{monochrome}: solo grises, para blanco y negro}, label={fig:secuencia-monochrome}, pos=t]
  sequenceDiagram
  autonumber
  actor Usuario
  participant texforge
  participant Validador
  participant Motor
  Usuario->>texforge: texforge build
  activate texforge
  texforge->>Validador: analiza el fuente
  activate Validador
  Validador-->>texforge: fuente válido
  deactivate Validador
  alt hay diagramas
  texforge->>Motor: prepara el motor
  activate Motor
  loop cada diagrama
  texforge->>Motor: renderiza diagrama
  Motor-->>texforge: figura lista
  end
  Motor-->>texforge: motor libre
  deactivate Motor
  else sin diagramas
  Note over texforge,Motor: compila directo
  end
  texforge->>Motor: compone el PDF
  activate Motor
  Motor-->>texforge: PDF final
  deactivate Motor
  texforge-->>Usuario: build/ listo
  deactivate texforge
\end{mermaid}

\begin{mermaid}[style=technical, width=0.62\linewidth, caption={\texttt{technical}: trazo uniforme, sin rellenos}, label={fig:secuencia-technical}, pos=t]
  sequenceDiagram
  autonumber
  actor Usuario
  participant texforge
  participant Validador
  participant Motor
  Usuario->>texforge: texforge build
  activate texforge
  texforge->>Validador: analiza el fuente
  activate Validador
  Validador-->>texforge: fuente válido
  deactivate Validador
  alt hay diagramas
  texforge->>Motor: prepara el motor
  activate Motor
  loop cada diagrama
  texforge->>Motor: renderiza diagrama
  Motor-->>texforge: figura lista
  end
  Motor-->>texforge: motor libre
  deactivate Motor
  else sin diagramas
  Note over texforge,Motor: compila directo
  end
  texforge->>Motor: compone el PDF
  activate Motor
  Motor-->>texforge: PDF final
  deactivate Motor
  texforge-->>Usuario: build/ listo
  deactivate texforge
\end{mermaid}

\begin{mermaid}[style=default, width=0.62\linewidth, caption={\texttt{default}: el aspecto nativo del motor}, label={fig:secuencia-default}, pos=t]
  sequenceDiagram
  autonumber
  actor Usuario
  participant texforge
  participant Validador
  participant Motor
  Usuario->>texforge: texforge build
  activate texforge
  texforge->>Validador: analiza el fuente
  activate Validador
  Validador-->>texforge: fuente válido
  deactivate Validador
  alt hay diagramas
  texforge->>Motor: prepara el motor
  activate Motor
  loop cada diagrama
  texforge->>Motor: renderiza diagrama
  Motor-->>texforge: figura lista
  end
  Motor-->>texforge: motor libre
  deactivate Motor
  else sin diagramas
  Note over texforge,Motor: compila directo
  end
  texforge->>Motor: compone el PDF
  activate Motor
  Motor-->>texforge: PDF final
  deactivate Motor
  texforge-->>Usuario: build/ listo
  deactivate texforge
\end{mermaid}

El preajuste \texttt{editorial} sirve para documentos que se leen en pantalla
o se imprimen a color y buscan un acabado sobrio; \texttt{monochrome} dibuja
solo en grises y es el indicado cuando el documento se imprime en blanco y
negro; \texttt{technical} viste los diagramas como planos o esquemas, sin
rellenos y con etiquetas de ancho fijo. El valor \texttt{default} conserva el
aspecto nativo de cada motor para no cambiar ningún documento existente. El
preajuste se fija una vez en \texttt{project.toml}, bajo \texttt{[diagrams]}
con \texttt{style = "..."}, y cada entorno puede cambiarlo con su propio
atributo \texttt{style}; omitir ambos equivale a \texttt{default}. Este
documento usa \texttt{editorial} en todo su recorrido.
\FloatBarrier

\subsection{Graphviz}

La figura~\ref{fig:arquitectura} dibuja la arquitectura como un grafo
dirigido con Graphviz~\cite{graphviz}, organizada en tres etapas: la interfaz
de línea de comandos, el preprocesador con sus tres motores y la salida que
produce el PDF. El motor de composición es Rust puro: no dibuja cajas de
\texttt{subgraph cluster}, así que las etapas se distinguen por el relleno de
sus nodos y la leyenda, no por un recuadro. Cada nodo nombra una pieza y las
flechas marcan dependencias.

\begin{graphviz}[width=1.0\linewidth, caption={Arquitectura de texforge: interfaz, preprocesador y salida}, label={fig:arquitectura}, pos=H]
  digraph Arquitectura {
    rankdir=LR
    node [style=filled, fillcolor=linen, color=sienna, fontsize=14]
    edge [color=sienna, fontsize=12]

    subgraph front {
      CLI [label="texforge CLI", shape=box, fillcolor=antiquewhite]
    }

    subgraph pre {
      Mermaid [label="Mermaid", fillcolor=cornsilk]
      Graphviz [label="Graphviz", fillcolor=cornsilk]
      D2 [label="D2", fillcolor=cornsilk]
      { rank=same; Mermaid; Graphviz; D2; }
    }

    subgraph back {
      Tectonic [label="tectonic", shape=box, fillcolor=darkseagreen]
      PDF [label="PDF", shape=doublecircle, fillcolor=white]
      { rank=same; Tectonic; PDF; }
    }

    CLI -> Mermaid [label="fuente", penwidth=2]
    CLI -> Graphviz [label="fuente", penwidth=2]
    CLI -> D2 [label="fuente", penwidth=2]
    Mermaid -> Tectonic [style=dashed, color=firebrick]
    Graphviz -> Tectonic [style=dashed, color=firebrick]
    D2 -> Tectonic [style=dashed, color=firebrick]
    Tectonic -> PDF [penwidth=2, color=firebrick, label="genera"]

    subgraph legend {
      E1 [label="origen: CLI", shape=box, fillcolor=antiquewhite]
      E2 [label="motor", shape=box, fillcolor=cornsilk]
      E3 [label="salida", shape=box, fillcolor=darkseagreen]
      E1 -> E2 [label="flujo", penwidth=2]
      E2 -> E3 [label="genera", style=dashed, color=firebrick]
    }
  }
\end{graphviz}
\FloatBarrier

\subsection{D2}

La figura~\ref{fig:pipeline} resume el trabajo como una tubería con la
sintaxis de D2: escribir, revisar, componer e inspeccionar, cada etapa con
su comando dentro de la caja, y el artefacto que produce la salida. La
arista a trazos marca el paso de verificación.

\begin{d2}[width=1.0\linewidth, caption={Tubería de trabajo, de la escritura a la inspección}, label={fig:pipeline}, pos=H]
  direction: down

  classes: {
    etapa: {
      style.border-radius: 4
    }
    control: {
      style.stroke-dash: 3
    }
  }

  autor: Autor {
    shape: person
    label: "Autor"
    class: etapa
  }

  flujo: Tubería {
    grid-columns: 2
    escribir: "Escribir\nfuente .tex" {
      shape: document
      class: etapa
    }
    revisar: "Revisar\ntexforge check" {
      shape: rectangle
      class: etapa
    }
    componer: "Componer\ntexforge build" {
      shape: hexagon
      class: etapa
    }
    inspeccionar: "Inspeccionar\ntexforge pdf info" {
      shape: rectangle
      class: etapa
    }
  }

  salida: Resultado {
    pdf: "build/\nPDF final" {
      shape: page
      class: etapa
    }
  }

  autor -> flujo.escribir: redacta
  flujo.escribir -> flujo.revisar {
    class: control
  }
  flujo.revisar -> flujo.componer
  flujo.componer -> flujo.inspeccionar
  flujo.inspeccionar -> salida.pdf: entrega

  leyenda: Leyenda {
    a: paso
    b: verificación
    a -> b: acción
    b -> a: "control" {
      class: control
    }
  }
\end{d2}
\FloatBarrier

\begin{notebox}
  Los tres entornos aceptan \texttt{label=} con o sin \texttt{caption}, y
  \texttt{texforge check} resuelve \texttt{\textbackslash ref} contra esas
  etiquetas. Si una figura pierde su etiqueta, el error aparece antes de
  compilar, no como una referencia sin resolver en el PDF.
\end{notebox}

\subsection{TikZ}

Para lo que ningún motor describe, el TikZ de siempre sigue disponible. La
figura~\ref{fig:pagina} dibuja la curva $f(x)=e^{-x^{2}/2}$ con su tangente
en $x=1$: el área bajo la curva queda sombreada, el punto de tangencia y los
ejes van anotados, y la paleta del documento (\texttt{docaccent} y los tonos
\texttt{bench}) mantiene el mismo acabado que el resto del capítulo. El
código fuente se muestra en el listado~\ref{lst:tikz-curva} para que pueda
copiarse tal cual.

\begin{figure}[H]
  \centering
  \begin{tikzpicture}
    \begin{axis}[
        width=0.82\linewidth,
        height=4.9cm,
        axis lines=middle,
        axis line style={benchsepia},
        xlabel={$x$},
        ylabel={$f(x)$},
        xmin=-2.6, xmax=2.6,
        ymin=-0.15, ymax=1.25,
        xtick={-2,-1,0,1,2},
        ytick={0.5,1},
        tick label style={font=\small\sffamily, color=docink},
        label style={font=\small\sffamily, color=docink},
        samples=120,
        domain=-2.5:2.5,
      ]
      \addplot[draw=none, fill=benchlamp, domain=-2.5:2.5] {exp(-x^2/2)} \closedcycle;
      \addplot[thick, color=docaccent] {exp(-x^2/2)};
      \addplot[dashed, color=benchsepia, domain=-0.2:2.2] {exp(-0.5) - exp(-0.5)*(x-1)};
      \addplot[only marks, mark=*, mark size=2pt, color=benchink] coordinates {(1,{exp(-0.5)})};
      \node[anchor=south west, font=\small\sffamily, color=docink] at (axis cs:1.05,0.62) {tangente en $x=1$};
      \node[anchor=south, font=\small\sffamily, color=benchsepia] at (axis cs:1.35,0.95) {$f(x)=e^{-x^{2}/2}$};
    \end{axis}
  \end{tikzpicture}
  \caption{Curva gaussiana con sombreado y tangente dibujada con TikZ}
  \label{fig:pagina}
\end{figure}

\begin{code}[lang=latex, numbers=true, caption={Código fuente de la figura~\ref{fig:pagina}}, label={lst:tikz-curva}]
\begin{tikzpicture}
  \begin{axis}[
    width=0.82\linewidth, height=4.9cm,
    axis lines=middle, axis line style={benchsepia},
    xlabel={$x$}, ylabel={$f(x)$},
    xmin=-2.6, xmax=2.6, ymin=-0.15, ymax=1.25,
    xtick={-2,-1,0,1,2}, ytick={0.5,1},
    tick label style={font=\small\sffamily, color=docink},
    label style={font=\small\sffamily, color=docink},
    samples=120, domain=-2.5:2.5,
  ]
    \addplot[draw=none, fill=benchlamp,
      domain=-2.5:2.5] {exp(-x^2/2)} \closedcycle;
    \addplot[thick, color=docaccent] {exp(-x^2/2)};
    \addplot[dashed, color=benchsepia,
      domain=-0.2:2.2]
      {exp(-0.5) - exp(-0.5)*(x-1)};
    \addplot[only marks, mark=*, mark size=2pt,
      color=benchink] coordinates {(1,{exp(-0.5)})};
    \node[anchor=south west, font=\small\sffamily,
      color=docink] at (axis cs:1.05,0.62)
      {tangente en $x=1$};
    \node[anchor=south, font=\small\sffamily,
      color=benchsepia] at (axis cs:1.35,0.95)
      {$f(x)=e^{-x^{2}/2}$};
  \end{axis}
\end{tikzpicture}
\end{code}

\subsection{Los preajustes}

El valor de \texttt{style} se fija una vez en \texttt{project.toml}, bajo
\texttt{[diagrams]}, y cada entorno puede cambiarlo con su propio atributo.
Los cuatro preajustes comparten el mismo criterio en Mermaid, Graphviz y D2:
\texttt{editorial} usa una paleta sobria con un solo acento para documentos a
color, \texttt{monochrome} dibuja solo en grises para la impresión en blanco y
negro, \texttt{technical} elimina rellenos y usa etiquetas de trazo uniforme
para esquemas, y \texttt{default} conserva el aspecto nativo de cada motor.
Omitir el atributo y la clave de proyecto equivale a \texttt{default}, de
modo que ningún documento existente cambia de aspecto.
